% Generated draft: requires mathematical and source review.
\documentclass[11pt]{article}
\usepackage[margin=25mm]{geometry}
\usepackage{fontspec}
\setmainfont{Latin Modern Roman}
\usepackage{amsmath,amssymb,mathtools,unicode-math}
\setmathfont{Latin Modern Math}
\usepackage{enumitem,needspace,xurl,hyperref}
\setlength{\parindent}{0pt}
\setlength{\parskip}{5pt}
\setlength{\emergencystretch}{3em}
\newcommand{\N}{\mathbb{N}}
\newcommand{\Z}{\mathbb{Z}}
\newcommand{\Q}{\mathbb{Q}}
\newcommand{\R}{\mathbb{R}}
\newcommand{\C}{\mathbb{C}}
\newcommand{\F}{\mathbb{F}}
\newcommand{\A}{\mathbb{A}}
\newcommand{\PP}{\mathbb{P}}
\newcommand{\E}{\mathbb{E}}
\newcommand{\eps}{\varepsilon}
\newcommand{\dd}{\,\mathrm{d}}
\begin{document}
\section*{UnsolvedMath ID: 800624}
\textbf{Source problem code:} catalogue entry
\subsection{Definitions and mathematical statement}
Let \(\mathcal M\) be the ideal of meager subsets of \(\R\), where \(\R\) has its usual topology. A set is meager if it is a countable union of nowhere dense sets; a set is nowhere dense if its closure has empty interior. Define
\[
\operatorname{cov}(\mathcal M)=\min\left\{|\mathcal A|:\mathcal A\subseteq\mathcal M\text{ and }\bigcup\mathcal A=\R\right\},
\qquad \mathfrak c=|\R|=2^{\aleph_0}.
\]
Let \(\omega_1\) be the least uncountable ordinal, and let \(\aleph_{\omega_1}\) denote the corresponding member of the aleph sequence. A \(\mathfrak c\)-Luzin set is a set \(L\subseteq\R\) such that
\[
|L|=\mathfrak c
\quad\text{and}\quad
(\forall M\in\mathcal M)\ |L\cap M|<\mathfrak c.
\]
The question is whether it is consistent that
\[
\operatorname{cov}(\mathcal M)=\aleph_1<\mathfrak c=\aleph_{\omega_1}
\quad\text{and there exists a }\mathfrak c\text{-Luzin set}.
\]
Under the conventional interpretation with ZFC as the ambient theory, this asks whether there is a model of ZFC satisfying both assertions, with all sets, cardinals, and quantifiers interpreted within that model. The supplied record does not explicitly specify the ambient theory or an intended relative-consistency hypothesis. The intersection bound is required separately for every meager set; no single cardinal below \(\mathfrak c\) is stipulated as a uniform bound.
\subsection{Short English statement}
Can the covering number of the meager ideal equal \(\aleph_1\), while the continuum equals \(\aleph_{\omega_1}\), and there is a set of continuum many reals meeting every meager set in fewer than continuum many points? The ambient consistency theory is not explicitly specified in the supplied source record.
\subsection{Sources}
\begin{itemize}
\item Supplied local UnsolvedMath registry record, permanent ID \(800624\), source problem code \texttt{catalogue entry}. Registry title: ``Selection Principles in Mathematics --- Problem 53 ([BarJu93]).'' The record attributes the problem to Tsaban, with proposed year 2006, and labels its status as open; this status is supplied curation, not an independently verified mathematical assessment.
\item Tsaban, \emph{Selection Principles and Special Sets of Reals: Open Problems} (2006), Problem 53. Supplied source URL: \url{https://arxiv.org/abs/math/0609601}. The record reports access on 2026-10-16 and extraction by \texttt{arxiv\_latex}; these provenance statements are supplied and have not been independently verified here.
\item The supplied background gives attribution/context as [BarJu93]. No full bibliographic entry for this reference is supplied, so its authorship, title, and precise role remain unresolved.
\end{itemize}
\end{document}
